% Source fingerprint: 084b160eef296743feca0908688a82cf7633fa248d7a0e5a5d68fbec6f1876ea
% T10: raw TeX cells preserved; no numeric highlighting.
\begin{table}[htbp]
\centering\small\renewcommand{\arraystretch}{1.18}\setlength{\tabcolsep}{4pt}
\caption[三种 Jacobi 因子的维数与用途]{三种 Jacobi 因子的维数与用途。线性公式在可微点用于导数矩阵。面积与余面积的源、目标维数相反，不能只交换转置符号而忽略积分对象。}
\label{b1:tab:visual-t10}
\begin{tabularx}{\linewidth}{@{}>{\raggedright\arraybackslash}p{3.1cm}>{\raggedright\arraybackslash}p{4.7cm}>{\raggedright\arraybackslash}X@{}}
\toprule
场景 & 非负因子 & 积分中的作用 \\
\midrule
\(T:\mathbb R^n\to\mathbb R^n\) & \(J_n(T)=|\det T|\) & 方阵换元的体积伸缩；与带符号行列式分开 \\
面积：\(T:\mathbb R^k\to\mathbb R^n\)，\(k\le n\) & \(J_k(T)=\sqrt{\det(T^{\mathsf T}T)}\) & 参数域的 \(k\) 维体积映到像上的 \(k\) 维面积；重叠另乘重数 \\
余面积：\(T:\mathbb R^n\to\mathbb R^m\)，\(m\le n\) & \(J_m(T)=\sqrt{\det(TT^{\mathsf T})}\) & 沿 \((n-m)\) 维纤维积分，再对目标变量积分 \\
秩亏 & 秩小于所需维数时相应因子为零 & 零因子结论不表示对应纤维为空 \\
\bottomrule
\end{tabularx}
\end{table}
